\documentclass[14pt,a4paper]{extarticle}
\usepackage[left=3cm,right=2cm,top=2cm,bottom=2cm,headheight=18pt]{geometry}
\usepackage{fontspec}
\setmainfont{Liberation Serif}
\setsansfont{Liberation Serif}
\setmonofont{Liberation Mono}
\usepackage{unicode-math}
\setmathfont{texgyretermes-math.otf}
\usepackage{amsmath,graphicx,booktabs,longtable,array,enumitem}
\usepackage{titlesec,fancyhdr}
\usepackage[hidelinks,unicode]{hyperref}
\usepackage{url}
\urlstyle{same}
\setlength{\parindent}{1.27cm}
\setlength{\parskip}{0pt}
\linespread{1}
\clubpenalty=10000
\widowpenalty=10000
\titleformat{\section}{\normalfont\bfseries}{\thesection.}{0.5em}{}
\titleformat{\subsection}{\normalfont\bfseries}{\thesubsection.}{0.5em}{}
\titlespacing*{\section}{0pt}{1.2em}{0.5em}
\titlespacing*{\subsection}{0pt}{0.9em}{0.3em}
\setlist{nosep,leftmargin=1.27cm}
\pagestyle{fancy}\fancyhf{}
\renewcommand{\headrulewidth}{0pt}
\renewcommand{\footrulewidth}{0pt}
\fancyhead[C]{\ifnum\value{page}>1\thepage\fi}
% Liberation Serif is a disclosed substitute, not Times New Roman.
% Drafting geometry and pagination: O‘RQ-682 appended methodology §§43–51.
\hyphenpenalty=10000
\exhyphenpenalty=10000
\pretolerance=10000
\tolerance=10000
\setlength{\emergencystretch}{4em}
\renewcommand{\normalsize}{\fontsize{14bp}{16.8bp}\selectfont}
\normalsize
\begin{document}
\noindent Author: Abduxoliq Akmal ugli Ashuraliyev. Edition of 1 October 2026.

\begin{flushright}Draft translation\end{flushright}
\begin{center}\textbf{PROPOSAL FOR INCLUSION\\IN A DRAFT NORMATIVE LEGAL ACT}\par
\medskip\textbf{On the procedure for verifying calculations of multidimensional family assessment}\end{center}
Author’s proposal for inclusion in the act being prepared under paragraph 14 of PF-193. English companion translation of the primary Russian submission. The Uzbek companion is supplied for consideration in the official state-language drafting process. The developer determines the final type of act and approval particulars. The provisions below are proposed and are not presented as existing law.

\noindent\hspace*{1.27cm}1. Approve, in accordance with the annex, the procedure for verifying calculations of the multidimensional assessment methodology used to include families in the Social Registry, assign categories and refer them to relevant services.\par
\noindent\hspace*{1.27cm}2. When preparing the multidimensional assessment methodology for approval in the prescribed manner, the National Agency for Social Protection under the President of the Republic of Uzbekistan (hereinafter, the Agency), jointly with the Ministry of Economy and Finance of the Republic of Uzbekistan (hereinafter, the Ministry), shall set out in a single description the criteria, calculation rules, category boundaries and their legal and scientific grounds. New criteria or grounds for refusal shall not be introduced independently through technical software settings.\par
\noindent\hspace*{1.27cm}3. The National Agency for Social Protection shall ensure verification of the conformity of calculations with the approved methodology, the ability to reproduce results and preservation of the grounds for decisions. Verification shall not constitute grounds for extending statutory periods for considering families’ applications.\par
\noindent\hspace*{1.27cm}4. The Agency, jointly with the Ministry, shall organize preparation of the methodology and verification materials under Article 22 of the Law “On Normative Legal Acts”. Consenting citizens with the necessary knowledge and experience may participate in the drafting group; scientific recommendations and specialist advice may be obtained. The Agency shall assign methodological, system and legal responsibilities and approve an internal technical regulation for criterion passports, source mapping, test inventories, protocols and defect records. That regulation shall implement this Procedure and shall not introduce eligibility criteria, refusal grounds, applicant obligations or changes to statutory periods. Drafting participation grants no authority to approve criteria or obtain protected records.\par
\noindent\hspace*{1.27cm}5. The Agency and Ministry shall reconcile the calculation description with registry assessment, category assignment, assistance, complaints and information exchange in the principal draft before the affected rule is first applied. Inconsistency in a higher-ranking act shall be submitted to its competent adopting body. Technical repair shall implement the rule legally applicable on the decision date; it shall not extend superseded criteria, change commencement dates or confer an employee power to enter a manual category decision outside the lawful procedure.\par
\noindent\hspace*{1.27cm}6. Verification shall be included in the existing system-development, technical-approval and change-management process. Existing specifications, tests, records and statutory oversight shall be reused where they provide equivalent evidence; duplicate controls and procurements shall not be required. The Agency and Ministry shall identify the incremental work, assigned staff time, contract coverage and lawful financing before committing additional expenditure. Any proposed processing change shall be supported by the relevant workload and timing records and, where necessary, a capacity test within the authorized environment. A costing formula, staff establishment or aggregate budget shall not substitute for that evidence. This Procedure establishes no separate nationwide manual processing service, parallel information system, new applicant charge or additional application stage.\par
\noindent\hspace*{1.27cm}7. Before consequential use, the Agency shall document the lawful rule, required technical approvals, control results and readiness of the existing explanation and objection routes in its release record under paragraph 25. A failed control shall be addressed under paragraph 26. No empirical welfare study or separate public verification report is a commencement condition. Technical acceptance shall not replace approvals required by legislation, change the 1 January 2027 commencement or extend applicant periods.\par
\clearpage\setcounter{page}{1}
\begin{flushright}\fontsize{12bp}{14.4bp}\selectfont Annex to Material 1\\Author’s proposal\\Edition of 1 October 2026\end{flushright}
\begin{center}\textbf{Procedure for verifying multidimensional assessment calculations}\end{center}
\subsection*{Chapter 1. General provisions}
\noindent\hspace*{1.27cm}1. This Procedure concerns conformity of calculations to the approved methodology, traceability of the data used and preservation of decision grounds. Control cases are checked before deployment and relevant changes within the existing release process, not by imposing a second manual calculation of every application. Sensitivity or welfare studies require a separately agreed question, lawful inputs and resources; they are not additional conditions for receiving assistance or launching the approved rule. Service-referral outputs are checked insofar as they follow from that rule; sector-specific service conditions remain.\par
\noindent\hspace*{1.27cm}2. This Procedure uses the following shortened names: the “Yagona milliy ijtimoiy himoya” information system (hereinafter, YAMIH) and the Agency’s “Inson” social services centres (hereinafter, Inson centres). The following definitions apply:\par
\textbf{methodology description} means a document recording assessment criteria, data sources, calculation rules, conditions for category assignment and periods of application;

\textbf{calculation version} means a linked record of the approved methodology identifier, the software release identifier and the configuration of effective-dated parameters used in a calculation;

\textbf{verification protocol} means a record of verification conditions, results, identified discrepancies and measures taken;

\textbf{unresolved data} means data not received, not confirmed, mutually contradictory or failing the methodology’s updating requirements;

\textbf{control calculation} means the result of independent recalculation under the approved rules.

\noindent\hspace*{1.27cm}3. The calculation methodology shall not replace the authority to award assistance or services to a family. Category, entitlement to particular assistance and its amount shall be determined separately under legislation.\par
\subsection*{Chapter 2. Methodology description and calculation versions}
\noindent\hspace*{1.27cm}4. For each criterion, the description shall identify its legal basis, purpose, authoritative source, unit, reference period, updating requirement and calculation rule. Applicable dates, variable parameters, duplicate or overlapping information and exceptions shall be determined consistently with the approved methodology. The Agency’s internal technical regulation shall prescribe passport fields, source mapping and precision checks; those forms shall not themselves establish a new rule or refusal ground.\par
\noindent\hspace*{1.27cm}5. The description shall cover formulas, conditional rules, coefficients, dependencies between criteria, exceptions, rounding and category boundaries. It shall specify the category applicable at exact equality. After lawful exceptions and precedence rules are applied, verification shall establish that a case is not assigned contradictory categories and that no unresolved gap remains between inclusion and refusal conditions. A case requiring clarification shall not be confused with an inclusion or refusal decision. Software shall not apply an assessment rule absent from the description.\par
\noindent\hspace*{1.27cm}6. For every version, the approval basis, beginning and end of application, differences from the previous version and verification protocol shall be preserved. A change affecting legal consequences shall be approved by the competent authority in the prescribed manner. The legal methodology, software release and parameter configuration shall have separately identifiable records linked to the calculation version, with their legal effective dates and technical deployment dates distinguished. A software correction implementing an unchanged lawful rule shall not be described as a new policy criterion. The configuration shall identify source mappings, conversions and rounding; changes capable of affecting results shall be verified.\par
\noindent\hspace*{1.27cm}7. The electronic description shall identify the published act approving the methodology, its application dates and the calculation version. Disclosable rule information shall be accessible through the Agency’s existing official information resource; an authoritative published description may be linked rather than copied into a second report. Protected parts remain available only to authorized verification and review personnel. This paragraph requires no publication of source code, protected records or infrastructure/security details, and creates no independent criterion or exception. Lawful temporary outage rules are identified in the version record under paragraph 11.\par
\subsection*{Chapter 3. Data quality and outages}
\noindent\hspace*{1.27cm}8. Each calculation shall record the source, reference period, receipt time and verification status of the data used. A genuine zero shall be distinguished from unresolved data. The date or period described by a source record shall be distinguished from its receipt time. Necessary input values or an authorized preserved source snapshot shall permit reconstruction of the calculation, subject to paragraph 28. Confirmed values, data not received, updating required, conflicting data and invalid format are recorded separately. Each used record retains its source and record identifiers, relevant period, receipt date and applied transformation. Conflicts preserve both values and the basis for selection; later receipt alone does not establish greater reliability. A staff correction records its author, time and lawful basis.\par
\noindent\hspace*{1.27cm}9. Unresolved data shall not automatically be equated with a value that worsens the family’s position. Where a normative value expressly provided by legislation is applied, its legal basis and reason for application shall be shown separately in the calculation.\par
\noindent\hspace*{1.27cm}10. Where unresolved data prevent a calculation, an authorized Inson employee shall register the issue in YAMIH and refer it to the responsible source owner or operator. Software defects shall be referred to the system unit; uncertainty in a legal rule shall be examined by the methodology unit and legal service and, where official interpretation or amendment is necessary, submitted to the competent body. Agency staff shall not replace that body’s powers. The responsible officer shall track the applicable statutory period, inform the family of the result and review route, and escalate an unresolved matter before that period expires. The status is not a new refusal or termination ground. A family document replacing state-system data may be required only where legislation permits.\par
\noindent\hspace*{1.27cm}11. Where a temporary algorithm change provided by legislation is introduced because of disrupted information exchange or incomplete receipt of electronic data, an authorized official shall record its basis, scope, start time, temporary calculation rules and restoration conditions. Calculations using the temporary version shall be separately identified. After restoration, verification and necessary recalculation shall follow the procedure established by legislation. The temporary record links affected sources, calculations and versions and identifies the deciding officer, legal authority and termination date. On restoration, temporary and restored calculations are compared; differences and individual decisions of the competent authority are retained. Closure is evidenced by source availability and completion of required checks, not erasure of the temporary record.\par
\subsection*{Chapter 4. Verification of calculations}
\noindent\hspace*{1.27cm}12. Before deployment or a change capable of affecting the score, category or decision ground, the implementation shall be compared with independently justified control answers derived from the approved legal rule. The existing acceptance and technical-compliance tests shall incorporate applicable boundaries, exceptions, their interactions, data-quality states and effective dates. A documented equivalent test shall be reused rather than repeated solely for this Procedure. Expected answers shall not be obtained solely from the tested implementation. The internal technical regulation shall prescribe case fields and reconciliation of planned and executed cases. An omitted applicable path, empty applicable comparison set or unresolved rule conflict shall not be represented as a successful check.\par
\noindent\hspace*{1.27cm}13. The purpose, period and cases covered, selection rules, comparison criteria and result-assessment conditions shall be specified before results are obtained. Subsequent changes shall be recorded with reasons. Synthetic test results shall not be used as evidence that real families’ needs are correctly identified.\par
\noindent\hspace*{1.27cm}14. Score and category discrepancies shall be counted separately. Reporting shall state the number of cases changing category because of the data or methodology changes examined, the relevant total number of cases and the nature of the change. Unresolved cases shall not be concealed and shall be distinguished from cases with a determined category. The category-disagreement denominator shall contain only comparable cases with determinate categories in both calculations under the same approved rule and identical inputs. A sensitivity denominator shall contain only paired cases with determinate categories before and after the specified change. Each measure shall disclose its numerator, denominator, excluded and unresolved counts; a zero denominator shall be reported as undefined, not zero.\par
\noindent\hspace*{1.27cm}15. Measures of incorrect identification of need shall be calculated only where independently justified comparison data exist. A previous automated decision, old category or complaint outcome shall not itself constitute a correct assessment of need. Complaint outcomes shall not be interpreted as an error measure covering families that did not complain.\par
\noindent\hspace*{1.27cm}16. The protocol shall identify the approved rule and tested implementation, the scope and basis of checks, responsible persons, results, untested or unresolved matters, identified effects and corrective decisions. It shall distinguish conformity to the rule, data quality and completeness of decision grounds from empirical identification of need. Intermediate discrepancies shall remain visible even where totals coincide. Agreement between programs alone does not establish a lawful expected answer. The internal technical regulation shall prescribe evidence and record formats; protected inputs shall remain in the lawfully accessible part.\par
\noindent\hspace*{1.27cm}17. Where calculations contrary to the approved methodology are identified, the software version shall be corrected and retested. Affected cases shall be identified and reconsidered lawfully by the competent authority. An unverified version shall not be used as a version confirmed to conform; this rule shall not justify stopping application processing or lawful assistance. The defect record identifies the rule and version, reproducible case, affected decisions, responsible person, correction, applicable period and retest result. Affected decisions are selected by the versions, dates and conditions used, not solely the complaint revealing the defect. Changed components and dependent calculation paths are retested; if the impact domain cannot be determined, the whole applicable calculation chain is checked. Closure requires verification evidence and a record of the required individual-decision review.\par
\subsection*{Chapter 5. Grounds for decisions and reconsideration}
\noindent\hspace*{1.27cm}18. Together with a family’s result, the version, relevant criteria, data sources and periods, score, category rule, normative values and basis of any authorized employee correction shall be preserved. Recalculation shall not erase the previous record. The record shall retain necessary intermediate contributions, inclusion or exclusion reasons, conversion and rounding stages and the identifiers specified in paragraph 6. Offsetting component discrepancies shall not be concealed by an identical total.\par
\noindent\hspace*{1.27cm}19. Together with the decision result, the applicant shall receive its legal and factual grounds, version, principal values and periods, reasons for applying a normative value or exception, and objection procedure through the personal electronic account. A person not using that account shall be able to obtain the information through an Inson centre. Protected third-party data shall not be disclosed. An SMS result notification shall not replace this explanation and access. The explanation includes decision identifier and date, applicable legal rule, the applicant’s relevant values and periods, principal transformations, category and conditions for particular assistance or referral, competent objection recipient and statutory period. Imputed income is distinguished from confirmed actual income. For a decision based solely on automated processing, the legally permitted basis and procedure for explaining that processing and its possible legal consequences are recorded before use. A lawful restriction on particular details does not prevent provision of the remaining accessible grounds.\par
\noindent\hspace*{1.27cm}20. An objection to a decision based solely on automated processing shall be considered by the data owner or operator, who shall notify the subject of the result in writing within ten days. Such an objection may be routed through YAMIH or an Inson centre; system registration shall not delay the running of the period. A correction request shall be referred to the authoritative source. On the subject’s request, the owner or operator shall amend or supplement personal data no later than three days from the request; identified inaccuracies shall be corrected immediately. A council complaint shall be considered within the council’s statutory subject matter. A computational objection outside that competence shall not be left unconsidered for that reason but shall be referred to the competent authority. Lawful court and other review procedures and periods shall be preserved. This paragraph introduces no additional mandatory pre-court stage.\par
\subsection*{Chapter 6. Research, reporting and oversight}
\noindent\hspace*{1.27cm}21. Before data are used for research or scientific verification, the lawful processing basis, purpose, necessary volume and safeguards shall be specified. Research anonymisation shall comply with legislation. Replacing an identifier alone shall not constitute anonymisation. Special requirements for health data shall be observed. The research-authorization record specifies the lawful question, necessary fields and period, authorized operator, permitted actions, retention period and disclosure controls. A citizen’s drafting participation does not itself constitute a family’s consent to data processing. Health information is used only under special statutory conditions. Research calculation does not change a category or payment without a separate lawful individual decision.\par
\noindent\hspace*{1.27cm}22. Research calculations shall be performed in an Agency-controlled environment. This Procedure grants an external scientific participant no authority to obtain personal data. Protected data, small groups allowing re-identification and linked tables allowing re-identification shall not be published. Before supplied code is executed, the authorized operator checks its contents, version, required access and outbound connections. Processing uses the minimum necessary permissions; protected inputs must not be transmitted to external services or included in public logs. Execution authorization and disclosure-control results are recorded in the protocol. Supplying code does not authorize its automatic execution in the production system.\par
\noindent\hspace*{1.27cm}23. The verification protocol shall be maintained internally with the existing release and oversight records. It shall identify the rule and software version, date, scope, executed and excluded cases, comparable and unresolved counts, demonstrated discrepancies, corrective action and limitations. Equivalent existing records may be referenced without creating duplicate journals. Information shall be disclosed through existing lawful reporting, applicant-access and oversight procedures. This paragraph establishes no separate report before each release, new periodic statistical publication or new household-data collection. A public extract, where lawfully required or authorized, shall undergo disclosure review and shall identify its scope and limitations; protected or linkable family records shall be excluded.\par
\noindent\hspace*{1.27cm}24. Oversight of version retention, verification results and corrections shall be organized within existing statutory powers. Verification results shall inform proposals for improving the methodology; they shall not be an independent legal basis for automatically reducing assistance.\par
\subsection*{Chapter 7. Acceptance, correction and retention}
\noindent\hspace*{1.27cm}25. The Agency head or a lawfully authorized deputy shall confirm calculation conformity in the existing release record, supported by methodological, system and legal conclusions. The record shall link approvals required by legislation, the applicable legal dates, independent control results, reused equivalent checks and any incremental work and resources. It shall identify unresolved issues and the action under paragraph 26. This internal confirmation shall neither substitute for competent technical or cybersecurity approval nor authorize eligibility policy or official interpretation. A demonstrated decision-affecting discrepancy shall not be treated as an immaterial remaining task.\par
\noindent\hspace*{1.27cm}26. A calculation release with a demonstrated score, category or decision-ground discrepancy under the same approved rule, identical inputs and approved rounding shall be corrected and the affected paths retested before their deployment. Unaffected paths may continue only where their separation and conformity are evidenced. A previous software release may be retained or restored only if it implements the rule legally applicable to the decision and its relevant controls and operational suitability are verified; age of the release does not preserve an expired methodology. Where a defect is discovered in operation, the system unit shall record its scope and affected applications, prioritize repair and refer affected individual decisions for lawful review. Receipt dates, statutory periods and the competent decision procedure remain; a technical hold is not a new refusal ground or authority to stop lawful assistance. The information-exchange outage power applies only within its statutory subject matter. No manual category-entry power, use of superseded assessment rules, debt exemption or suspension of the January commencement is conferred. An absent or conflicting legal rule shall be submitted by the Agency and Ministry to the competent adopting or interpreting body for resolution before its first application; an internal setting or interpretation shall not amend law. A technical contingency cannot solve the absence of a lawful rule. Finite controls verify the stated scope, not all inputs or empirical need accuracy.\par
\noindent\hspace*{1.27cm}27. New legal criteria shall apply from the date established by the competent act. An error in a previous decision shall be checked under the legal rules applicable to that decision; a new criterion or later-period parameter shall not automatically be applied retrospectively. Recalculation alone shall not be an independent basis for recovering a previously lawfully awarded payment. Changes, continuation or recovery shall be decided on the relevant legal basis and by an individual decision. The review record identifies whether source facts, implementation of the previous rule or the legal rule itself changed; it preserves previous and corrected components, the applicable date and grounds for amending the individual decision. The inventory of affected cases and their referral to the competent body are linked to the defect record under paragraph 17. Without lawful grounds to change a payment, a technical correction does not become a payment instruction.\par
\noindent\hspace*{1.27cm}28. The Agency shall determine record categories, lawful retention period and basis, access powers, correction history and anonymisation or destruction procedures. This Procedure does not authorize indefinite retention. When the retention basis or period expires, data shall be destroyed or, where lawfully permitted, anonymised in accordance with archival and personal-data requirements. Where destruction is mandatory, anonymisation shall not replace it. Records necessary for a complaint or lawful oversight shall be protected under current retention requirements. Retention categories separately cover inputs, calculation traces, control protocols, corrections and submissions. Each identifies purpose, legal basis of period, responsible person, permitted access roles and termination method. Access to and modification of protected records are logged; checkers receive the necessary volume within their powers. A control copy does not authorize retaining personal information after its lawful basis ends.\par
\clearpage
\section*{Compact alternative for incorporation in the principal draft}
Recommended minimum for normative incorporation into the act mandated by PF-193 paragraph 14. The five paragraphs below replace the full seven operative provisions and 28-paragraph annex; they are not additional cumulative requirements. Detailed passports, cases and protocols belong in internal technical documentation. Where a specific requirement is already secured by an applicable equivalent norm of the required legal force, consistent with higher-ranking acts, the developer retains that norm and identifies it in the principal draft or its justification; repeating its text is unnecessary. Final numbering follows the place of incorporation.

1. When preparing the multidimensional methodology for approval in the prescribed manner, the National Agency for Social Protection under the President of the Republic of Uzbekistan (hereinafter, the Agency) and the Ministry of Economy and Finance of the Republic of Uzbekistan (hereinafter, the Ministry) shall ensure legal determination of the criteria and conditions for assigning families to categories. That determination covers category-affecting thresholds, equality, exception precedence, units and periods, normative imputation, permissible inputs and rounding insofar as they establish an inclusion or refusal condition. It identifies the competent instrument, application date and references to retained current rules. Software descriptions and technical acceptance shall neither supply a missing legal condition nor amend a higher-ranking norm. Normative provisions shall be approved and officially published in the prescribed manner; an applicable state-registration requirement shall be observed for a departmental normative act. Express statutory powers to change the algorithm temporarily remain within their subject matter.

2. The Agency shall ensure conformity to the lawfully applicable methodology and check affected rules before first use or a relevant software change through its existing release process. Equivalent existing checks shall be reused without duplication. Results shall link the legal rule, dates, tested release and independently justified expected answer. A demonstrated consequential discrepancy shall be corrected and retested; another build may be used only where it conforms to the rule lawfully applicable on the decision date and its operational suitability is confirmed. Affected individual decisions shall be referred for lawful reconsideration by the competent body. Required technical approvals, statutory periods, commencement date and assistance conditions remain.

3. With each consequential result, the Agency shall preserve the lawfully necessary record of the applied legal rule, calculation version, input values and periods, principal transformations, imputation, exception and category grounds sufficient to reconstruct the calculation. Through the existing personal account or the Agency’s Inson social services centre, the applicant shall receive accessible legal and factual grounds, relevant personal values and periods, and the objection route; an SMS result shall not replace explanation. Statutory data-correction and objection periods remain. Protected third-party information shall not be disclosed; retention, access and termination of retention shall comply with legislation.

4. A change to a legal assessment condition shall be submitted to the competent body and applied from its lawful date; software repair implements the unchanged lawful rule. Where legislation expressly permits temporary algorithm changes during technical faults, exchange interruptions or incomplete receipt of electronic data, an authorized official shall record the basis, affected sources and cases, temporary rule, beginning and termination conditions and, after restoration, necessary checks and lawful review. This proposal does not enlarge that power. A missing or contradictory legal norm shall be resolved by the competent adopting or interpreting body before the affected condition is first applied; interpretation shall not insert amendments. Technical waiting creates no refusal ground, period extension, authority to apply repealed criteria or debt exemption.

5. The Agency shall organize checks, forms and records within its existing powers and development, approval and oversight processes. Internal technical documents shall establish no new applicant obligations or legal criteria; equivalent specifications, tests, records and channels shall be reused. Before committing additional expenditure, the Agency and Ministry shall identify only uncovered work, assigned staff time, contractual coverage and lawful financing; a processing change shall be supported by its relevant actual workload and timing evidence. No separate manual service, additional application stage or new mandatory public statistical programme is established.

\end{document}
